\chapter{Предел}

% -----------------------------------------------------------------------------
\lecture{Предел последовательности}

\subsection{Определения и примеры}

\begin{definition}{Последовательность}{c3-sequence}
  Функция $f\colon\N\to X$, областью определения которой является множество
  натуральных чисел, называется последовательностью.

  Значения $f(n)$ функции $f$ называются членами последовательности. Их
  принято обозначать символом элемента того множества, в которое идет
  отображение, наделяя символ соответствующим индексом аргумента,
  $x_n:=f(n)$. Саму последовательность в связи с этим обозначают символом
  $\{x_n\}$, а также записывают в виде
  \[
    x_1,x_2,\ldots,x_n,\ldots
  \]
  и называют последовательностью в $X$ или последовательностью элементов
  множества $X$. Элемент $x_n$ называется $n$-м членом последовательности.
\end{definition}

\begin{definition}{Предел числовой последовательности}{c3-sequence-limit}
  Число $A\in\R$ называется пределом числовой последовательности
  $\{x_n\}$, если для любой окрестности $V(A)$ точки $A$ существует такой
  номер $N$ (выбираемый в зависимости от $V(A)$), что все члены
  последовательности, номера которых больше $N$, содержатся в указанной
  окрестности точки $A$:
  \[
    \lim_{n\to\infty}x_n=A
    \;:=\;
    \forall V(A)\ \exists N\in\N\ \forall n>N\;(x_n\in V(A)).
  \]
  Равносильно: число $A\in\R$ называется пределом последовательности
  $\{x_n\}$, если для любого $\varepsilon>0$ существует номер $N$ такой,
  что при всех $n>N$ имеем $|x_n-A|<\varepsilon$:
  \[
    \lim_{n\to\infty}x_n=A
    \;:=\;
    \forall\varepsilon>0\ \exists N\in\N\ \forall n>N\
    (|x_n-A|<\varepsilon).
  \]
\end{definition}

\begin{definition}{Сходящаяся последовательность}{c3-convergent-sequence}
  Если $\lim\limits_{n\to\infty}x_n=A$, то говорят, что последовательность
  $\{x_n\}$ сходится к $A$ или стремится к $A$, и пишут $x_n\to A$ при
  $n\to\infty$.

  Последовательность, имеющая предел, называется сходящейся.
  Последовательность, не имеющая предела, называется расходящейся.
\end{definition}

\subsection{Свойства предела последовательности}

\begin{definition}{Финально постоянная последовательность}{c3-eventually-constant-sequence}
  Если существуют число $A$ и номер $N$ такие, что $x_n=A$ при любом
  $n>N$, то последовательность $\{x_n\}$ будем называть финально
  постоянной.
\end{definition}

\begin{definition}{Ограниченная последовательность}{c3-bounded-sequence}
  Последовательность $\{x_n\}$ называется ограниченной, если существует
  число $M$ такое, что $|x_n|<M$ при любом $n\in\N$.
\end{definition}

\begin{theorem}{Общие свойства предела последовательности}{c3-sequence-limit-properties}
  \begin{enumerate}[label=\alph*)]
    \item Финально постоянная последовательность сходится.
    \item Любая окрестность предела последовательности содержит все члены
      последовательности, за исключением конечного их числа.
    \item Последовательность не может иметь двух различных пределов.
    \item Сходящаяся последовательность ограничена.
  \end{enumerate}
\end{theorem}

\begin{definition}{Арифметические операции над последовательностями}{c3-sequence-arithmetic}
  Если $\{x_n\}$, $\{y_n\}$ — две числовые последовательности, то их
  суммой, произведением и частным называются соответственно
  последовательности
  \[
    \{x_n+y_n\},\qquad \{x_ny_n\},\qquad
    \left\{\frac{x_n}{y_n}\right\}.
  \]
  Частное определено лишь при $y_n\ne0$, $n\in\N$.
\end{definition}

\begin{theorem}{Арифметические свойства предела последовательности}{c3-sequence-limit-arithmetic}
  Пусть $\{x_n\}$, $\{y_n\}$ — числовые последовательности. Если
  \[
    \lim_{n\to\infty}x_n=A,
    \qquad
    \lim_{n\to\infty}y_n=B,
  \]
  то
  \begin{enumerate}[label=\alph*)]
    \item $\displaystyle\lim_{n\to\infty}(x_n+y_n)=A+B$;
    \item $\displaystyle\lim_{n\to\infty}x_ny_n=AB$;
    \item $\displaystyle\lim_{n\to\infty}\frac{x_n}{y_n}=\frac AB$,
      если $y_n\ne0$ $(n=1,2,\ldots)$ и $B\ne0$.
  \end{enumerate}
\end{theorem}

\begin{theorem}{Предельный переход и неравенства}{c3-sequence-limit-inequalities}
  \begin{enumerate}[label=\alph*)]
    \item Пусть $\{x_n\}$, $\{y_n\}$ — две сходящиеся
      последовательности, причем
      \[
        \lim_{n\to\infty}x_n=A,
        \qquad
        \lim_{n\to\infty}y_n=B.
      \]
      Если $A<B$, то найдется номер $N\in\N$ такой, что при любом $n>N$
      выполнено неравенство $x_n<y_n$.
    \item Пусть последовательности $\{x_n\}$, $\{y_n\}$, $\{z_n\}$ таковы,
      что при любом $n>N\in\N$ имеет место соотношение
      $x_n\le y_n\le z_n$. Если при этом последовательности $\{x_n\}$,
      $\{z_n\}$ сходятся к одному и тому же пределу, то последовательность
      $\{y_n\}$ также сходится и к этому же пределу.
  \end{enumerate}
\end{theorem}

\begin{corollary}{Предельный переход в неравенствах}{c3-sequence-inequality-corollary}
  Пусть $\lim\limits_{n\to\infty}x_n=A$ и
  $\lim\limits_{n\to\infty}y_n=B$. Если существует номер $N$ такой, что
  при любом $n>N$
  \begin{enumerate}[label=\alph*)]
    \item $x_n>y_n$, то $A\ge B$;
    \item $x_n\ge y_n$, то $A\ge B$;
    \item $x_n>B$, то $A\ge B$;
    \item $x_n\ge B$, то $A\ge B$.
  \end{enumerate}
\end{corollary}

\subsection{Вопросы существования предела последовательности}

\begin{definition}{Фундаментальная последовательность}{c3-cauchy-sequence}
  Последовательность $\{x_n\}$ называется фундаментальной (или
  последовательностью Коши), если для любого числа $\varepsilon>0$ найдется
  такой номер $N\in\N$, что из $n>N$ и $m>N$ следует
  $|x_m-x_n|<\varepsilon$.
\end{definition}

\begin{theorem}{Критерий Коши сходимости последовательности}{c3-cauchy-criterion-sequence}
  Числовая последовательность сходится тогда и только тогда, когда она
  фундаментальна.
\end{theorem}

\begin{definition}{Монотонные последовательности}{c3-monotone-sequences}
  Последовательность $\{x_n\}$ называется возрастающей, если
  $\forall n\in\N\;(x_n<x_{n+1})$; неубывающей, если
  $\forall n\in\N\;(x_n\le x_{n+1})$; невозрастающей, если
  $\forall n\in\N\;(x_n\ge x_{n+1})$; убывающей, если
  $\forall n\in\N\;(x_n>x_{n+1})$.

  Последовательности этих четырех типов называют монотонными
  последовательностями.
\end{definition}

\begin{definition}{Ограниченная сверху последовательность}{c3-sequence-bounded-above}
  Последовательность $\{x_n\}$ называется ограниченной сверху, если
  существует число $M$ такое, что $\forall n\in\N\;(x_n<M)$.
  Аналогично определяется последовательность, ограниченная снизу.
\end{definition}

\begin{theorem}{Вейерштрасса}{c3-weierstrass-monotone-sequence}
  Для того чтобы неубывающая последовательность имела предел, необходимо
  и достаточно, чтобы она была ограниченной сверху.
\end{theorem}

\begin{properties}{Невозрастающая последовательность}{c3-decreasing-sequence-limit}
  Для того чтобы невозрастающая последовательность имела предел, необходимо
  и достаточно, чтобы она была ограниченной снизу. В этом случае
  \[
    \lim_{n\to\infty}x_n=\inf_{n\in\N}x_n.
  \]
\end{properties}

\begin{corollary}{Предел корня из $n$}{c3-root-n-limit}
  \[
    \lim_{n\to\infty}\sqrt[n]{n}=1.
  \]
\end{corollary}

\begin{corollary}{Предел корня из положительного числа}{c3-root-a-limit}
  При любом $a>0$
  \[
    \lim_{n\to\infty}\sqrt[n]{a}=1.
  \]
\end{corollary}

\begin{definition}{Число $e$}{c3-number-e}
  \[
    e:=\lim_{n\to\infty}\left(1+\frac1n\right)^n.
  \]
\end{definition}

\begin{definition}{Подпоследовательность}{c3-subsequence}
  Если $x_1,x_2,\ldots,x_n,\ldots$ — некоторая последовательность, а
  $n_1<n_2<\ldots<n_k<\ldots$ — возрастающая последовательность
  натуральных чисел, то последовательность
  \[
    x_{n_1},x_{n_2},\ldots,x_{n_k},\ldots
  \]
  называется подпоследовательностью последовательности $\{x_n\}$.
\end{definition}

\begin{lemma}{Больцано--Вейерштрасса}{c3-bolzano-weierstrass}
  Каждая ограниченная последовательность действительных чисел содержит
  сходящуюся подпоследовательность.
\end{lemma}

\begin{definition}{Последовательность, стремящаяся к бесконечности}{c3-sequence-infinite-limit}
  Условимся писать $x_n\to+\infty$ и говорить, что последовательность
  $\{x_n\}$ стремится к плюс бесконечности, если для каждого числа $c$
  найдется номер $N\in\N$ такой, что $x_n>c$ при любом $n>N$.

  В логических обозначениях:
  \[
    (x_n\to+\infty):=\forall c\in\R\ \exists N\in\N\ \forall n>N\;(c<x_n),
  \]
  \[
    (x_n\to-\infty):=\forall c\in\R\ \exists N\in\N\ \forall n>N\;(x_n<c),
  \]
  \[
    (x_n\to\infty):=\forall c\in\R\ \exists N\in\N\ \forall n>N\;(c<|x_n|).
  \]
  В последних двух случаях говорят соответственно: последовательность
  $\{x_n\}$ стремится к минус бесконечности и последовательность $\{x_n\}$
  стремится к бесконечности.
\end{definition}

\begin{lemma}{О подпоследовательности}{c3-subsequence-lemma}
  Из каждой последовательности действительных чисел можно извлечь
  сходящуюся подпоследовательность или подпоследовательность, стремящуюся
  к бесконечности.
\end{lemma}

\begin{definition}{Нижний предел последовательности}{c3-liminf-sequence}
  Число
  \[
    l=\lim_{n\to\infty}\inf_{k\ge n}x_k
  \]
  называется нижним пределом последовательности $\{x_k\}$ и обозначается
  $\lim\limits_{k\to\infty}x_k$ с нижней чертой или
  $\liminf\limits_{k\to\infty}x_k$.

  С учетом конечных и бесконечных значений определение записывается так:
  \[
    \liminf_{k\to\infty}x_k:=
    \lim_{n\to\infty}\inf_{k\ge n}x_k.
  \]
\end{definition}

\begin{definition}{Верхний предел последовательности}{c3-limsup-sequence}
  \[
    \limsup_{k\to\infty}x_k:=
    \lim_{n\to\infty}\sup_{k\ge n}x_k.
  \]
\end{definition}

\begin{definition}{Частичный предел последовательности}{c3-partial-limit}
  Число (или символ $-\infty$ или $+\infty$) называют частичным пределом
  последовательности, если в ней есть подпоследовательность, стремящаяся к
  этому числу.
\end{definition}

\begin{proposition}{О нижнем и верхнем пределах ограниченной последовательности}{c3-extreme-partial-limits-bounded}
  Нижний и верхний пределы ограниченной последовательности являются
  соответственно наименьшим и наибольшим из ее частичных пределов.
\end{proposition}

\begin{proposition}{О нижнем и верхнем пределах произвольной последовательности}{c3-extreme-partial-limits-general}
  Для любой последовательности нижний предел есть наименьший из ее
  частичных пределов, а верхний предел последовательности — наибольший из
  ее частичных пределов.
\end{proposition}

\begin{corollary}{Критерий существования конечного или бесконечного предела}{c3-liminf-limsup-criterion}
  Последовательность имеет предел или стремится к минус или плюс
  бесконечности в том и только в том случае, когда нижний и верхний пределы
  последовательности совпадают.
\end{corollary}

\begin{corollary}{Сходимость подпоследовательностей}{c3-subsequence-convergence}
  Последовательность сходится тогда и только тогда, когда сходится любая ее
  подпоследовательность.
\end{corollary}

\begin{corollary}{Лемма Больцано--Вейерштрасса}{c3-bw-from-limits}
  Лемма Больцано--Вейерштрасса как в узкой, так и в расширенной
  формулировке вытекает из утверждения о нижнем и верхнем пределах
  ограниченной последовательности и утверждения о нижнем и верхнем
  пределах произвольной последовательности соответственно.
\end{corollary}

\subsection{Начальные сведения о рядах}

\begin{definition}{Ряд}{c3-series}
  Выражение
  \[
    a_1+a_2+\ldots+a_n+\ldots
  \]
  обозначают символом
  \[
    \sum_{n=1}^{\infty}a_n
  \]
  и обычно называют рядом или бесконечным рядом.
\end{definition}

\begin{definition}{Члены ряда}{c3-series-terms}
  Элементы последовательности $\{a_n\}$, рассматриваемые как элементы ряда,
  называют членами ряда; элемент $a_n$ называют $n$-м членом ряда.
\end{definition}

\begin{definition}{Частичная сумма ряда}{c3-partial-sum}
  Сумму
  \[
    s_n=\sum_{k=1}^{n}a_k
  \]
  называют частичной суммой ряда или, когда желают указать ее номер,
  $n$-й частичной суммой ряда.
\end{definition}

\begin{definition}{Сходящийся ряд}{c3-convergent-series}
  Если последовательность $\{s_n\}$ частичных сумм ряда сходится, то ряд
  называется сходящимся. Если последовательность $\{s_n\}$ не имеет
  предела, то ряд называют расходящимся.
\end{definition}

\begin{definition}{Сумма ряда}{c3-series-sum}
  Предел $\lim\limits_{n\to\infty}s_n=s$ последовательности частичных сумм,
  если он существует, называется суммой ряда. Именно в этом смысле
  понимается запись
  \[
    \sum_{n=1}^{\infty}a_n=s.
  \]
\end{definition}

\begin{theorem}{Критерий Коши сходимости ряда}{c3-cauchy-criterion-series}
  Ряд $a_1+\ldots+a_n+\ldots$ сходится тогда и только тогда, когда для
  любого $\varepsilon>0$ найдется такое число $N\in\N$, что из
  $m\ge n>N$ следует
  \[
    |a_n+\ldots+a_m|<\varepsilon.
  \]
\end{theorem}

\begin{corollary}{Изменение конечного числа членов ряда}{c3-finite-series-change}
  Если в ряде изменить только конечное число членов, то получающийся при
  этом новый ряд будет сходиться, если сходился исходный ряд, и будет
  расходиться, если исходный ряд расходился.
\end{corollary}

\begin{corollary}{Необходимое условие сходимости ряда}{c3-series-necessary-condition}
  Для того чтобы ряд $a_1+\ldots+a_n+\ldots$ сходился, необходимо, чтобы
  его члены стремились к нулю при $n\to\infty$, то есть необходимо
  \[
    \lim_{n\to\infty}a_n=0.
  \]
\end{corollary}

\begin{definition}{Абсолютная сходимость}{c3-absolute-convergence}
  Ряд $\sum\limits_{n=1}^{\infty}a_n$ называется абсолютно сходящимся,
  если сходится ряд $\sum\limits_{n=1}^{\infty}|a_n|$.
\end{definition}

\begin{properties}{Абсолютная сходимость и сходимость}{c3-absolute-implies-convergence}
  Если ряд сходится абсолютно, то он сходится.
\end{properties}

\begin{theorem}{Критерий сходимости рядов с неотрицательными членами}{c3-nonnegative-series-criterion}
  Ряд $a_1+\ldots+a_n+\ldots$, члены которого — неотрицательные числа,
  сходится тогда и только тогда, когда последовательность его частичных
  сумм ограничена сверху.
\end{theorem}

\begin{theorem}{Теорема сравнения}{c3-series-comparison}
  Пусть $\sum\limits_{n=1}^{\infty}a_n$,
  $\sum\limits_{n=1}^{\infty}b_n$ — два ряда с неотрицательными членами.
  Если существует номер $N\in\N$ такой, что при любом $n>N$ имеет место
  неравенство $a_n\le b_n$, то из сходимости ряда
  $\sum\limits_{n=1}^{\infty}b_n$ вытекает сходимость ряда
  $\sum\limits_{n=1}^{\infty}a_n$, а из расходимости ряда
  $\sum\limits_{n=1}^{\infty}a_n$ вытекает расходимость ряда
  $\sum\limits_{n=1}^{\infty}b_n$.
\end{theorem}

\begin{corollary}{Мажорантный признак Вейерштрасса абсолютной сходимости ряда}{c3-weierstrass-majorant}
  Пусть $\sum\limits_{n=1}^{\infty}a_n$ и
  $\sum\limits_{n=1}^{\infty}b_n$ — два ряда. Пусть существует номер
  $N\in\N$ такой, что при любом $n>N$ имеет место соотношение
  $|a_n|\le b_n$. При этих условиях для абсолютной сходимости ряда
  $\sum\limits_{n=1}^{\infty}a_n$ достаточно, чтобы ряд
  $\sum\limits_{n=1}^{\infty}b_n$ сходился.
\end{corollary}

\begin{corollary}{Признак Коши}{c3-cauchy-root-test}
  Пусть $\sum\limits_{n=1}^{\infty}a_n$ — данный ряд и
  \[
    \alpha=\limsup_{n\to\infty}\sqrt[n]{|a_n|}.
  \]
  Тогда справедливы следующие утверждения:
  \begin{enumerate}[label=\alph*)]
    \item если $\alpha<1$, то ряд $\sum\limits_{n=1}^{\infty}a_n$
      абсолютно сходится;
    \item если $\alpha>1$, то ряд $\sum\limits_{n=1}^{\infty}a_n$
      расходится;
    \item существуют как абсолютно сходящиеся, так и расходящиеся ряды,
      для которых $\alpha=1$.
  \end{enumerate}
\end{corollary}

\begin{corollary}{Признак Даламбера}{c3-dalembert-ratio-test}
  Пусть для ряда $\sum\limits_{n=1}^{\infty}a_n$ существует предел
  \[
    \lim_{n\to\infty}\left|\frac{a_{n+1}}{a_n}\right|=\alpha.
  \]
  Тогда справедливы следующие утверждения:
  \begin{enumerate}[label=\alph*)]
    \item если $\alpha<1$, то ряд $\sum\limits_{n=1}^{\infty}a_n$
      сходится абсолютно;
    \item если $\alpha>1$, то ряд $\sum\limits_{n=1}^{\infty}a_n$
      расходится;
    \item существуют как абсолютно сходящиеся, так и расходящиеся ряды,
      для которых $\alpha=1$.
  \end{enumerate}
\end{corollary}

\begin{proposition}{Признак конденсации Коши}{c3-cauchy-condensation}
  Если $a_1\ge a_2\ge\ldots\ge0$, то ряд
  \[
    \sum_{n=1}^{\infty}a_n
  \]
  сходится тогда и только тогда, когда сходится ряд
  \[
    \sum_{k=0}^{\infty}2^k a_{2^k}
    =a_1+2a_2+4a_4+8a_8+\ldots
  \]
\end{proposition}

\begin{corollary}{Сходимость обобщенного гармонического ряда}{c3-p-series}
  Ряд
  \[
    \sum_{n=1}^{\infty}\frac1{n^p}
  \]
  сходится при $p>1$ и расходится при $p\le1$.
\end{corollary}

% -----------------------------------------------------------------------------
\lecture{Предел функции}

\subsection{Определения и примеры}

\begin{definition}{Предел функции по Коши}{c3-function-limit-cauchy}
  Будем (следуя Коши) говорить, что функция $f\colon E\to\R$ стремится к
  $A$ при $x$, стремящемся к $a$, или что $A$ является пределом функции $f$
  при $x$, стремящемся к $a$, если для любого числа $\varepsilon>0$
  существует число $\delta>0$ такое, что для любой точки $x\in E$ такой,
  что $0<|x-a|<\delta$, выполнено соотношение $|f(x)-A|<\varepsilon$:
  \[
    \forall\varepsilon>0\ \exists\delta>0\ \forall x\in E\
    \bigl(0<|x-a|<\delta\Rightarrow |f(x)-A|<\varepsilon\bigr).
  \]
\end{definition}

\begin{definition}{Проколотая окрестность}{c3-punctured-neighborhood}
  Проколотой окрестностью точки называется окрестность точки, из которой
  исключена сама эта точка.

  Если $U(a)$ — обозначение окрестности точки $a$, то проколотую
  окрестность этой точки будем обозначать символом $\mathring U(a)$.
  Множества
  \[
    U_E(a):=E\cap U(a),
    \qquad
    \mathring U_E(a):=E\cap\mathring U(a)
  \]
  будем называть соответственно окрестностью и проколотой окрестностью
  точки $a$ в множестве $E$.
\end{definition}

\begin{definition}{Предел функции в терминах окрестностей}{c3-function-limit-neighborhoods}
  \[
    \lim_{E\ni x\to a}f(x)=A
    \;:=\;
    \forall V_{\R}(A)\ \exists\mathring U_E(a)\
    \bigl(f(\mathring U_E(a))\subset V_{\R}(A)\bigr).
  \]
  Итак, число $A$ называется пределом функции $f\colon E\to\R$ при $x$,
  стремящемся по множеству $E$ к точке $a$ (предельной для $E$), если для
  любой окрестности точки $A$ найдется проколотая окрестность точки $a$ в
  множестве $E$, образ которой при отображении $f\colon E\to\R$ содержится
  в заданной окрестности точки $A$.
\end{definition}

\begin{proposition}{Равносильность определений предела по Коши и по Гейне}{c3-heine-limit}
  Соотношение
  \[
    \lim_{E\ni x\to a}f(x)=A
  \]
  имеет место тогда и только тогда, когда для любой последовательности
  $\{x_n\}$ точек $x_n\in E\setminus\{a\}$, сходящейся к $a$,
  последовательность $\{f(x_n)\}$ сходится к $A$.
\end{proposition}

\subsection{Свойства предела функции}

\begin{definition}{Постоянная и финально постоянная функции}{c3-eventually-constant-function}
  Функцию $f\colon E\to\R$, принимающую только одно значение, будем, как и
  прежде, называть постоянной. Функция $f\colon E\to\R$ называется
  финально постоянной при $E\ni x\to a$, если она постоянна в некоторой
  проколотой окрестности $\mathring U_E(a)$ точки $a$, предельной для
  множества $E$.
\end{definition}

\begin{definition}{Ограниченные и финально ограниченные функции}{c3-bounded-function}
  Функция $f\colon E\to\R$ называется ограниченной, ограниченной сверху,
  ограниченной снизу, если найдется число $C\in\R$ такое, что для любого
  $x\in E$ выполнено соответственно
  \[
    |f(x)|<C,\qquad f(x)<C,\qquad C<f(x).
  \]
  В случае, если первое, второе или третье из этих соотношений выполнено
  лишь в некоторой проколотой окрестности $\mathring U_E(a)$ точки $a$,
  функция $f\colon E\to\R$ называется соответственно финально ограниченной
  при $E\ni x\to a$, финально ограниченной сверху при $E\ni x\to a$,
  финально ограниченной снизу при $E\ni x\to a$.
\end{definition}

\begin{theorem}{Общие свойства предела функции}{c3-function-limit-properties}
  \begin{enumerate}[label=\alph*)]
    \item Если $f\colon E\to\R$ при $E\ni x\to a$ есть финально постоянная
      $A$, то $\lim\limits_{E\ni x\to a}f(x)=A$.
    \item Если существует $\lim\limits_{E\ni x\to a}f(x)$, то
      $f\colon E\to\R$ финально ограничена при $E\ni x\to a$.
    \item Если
      \[
        \lim_{E\ni x\to a}f(x)=A_1,
        \qquad
        \lim_{E\ni x\to a}f(x)=A_2,
      \]
      то $A_1=A_2$.
  \end{enumerate}
\end{theorem}

\begin{definition}{Арифметические операции над функциями}{c3-function-arithmetic}
  Если две числовые функции $f\colon E\to\R$, $g\colon E\to\R$ имеют
  общую область определения $E$, то их суммой, произведением и частным
  называются соответственно функции, определенные на том же множестве
  следующими формулами:
  \[
    (f+g)(x):=f(x)+g(x),
    \qquad
    (f\cdot g)(x):=f(x)g(x),
  \]
  \[
    \left(\frac fg\right)(x):=\frac{f(x)}{g(x)},
    \qquad g(x)\ne0\quad(x\in E).
  \]
\end{definition}

\begin{theorem}{Арифметические свойства предела функции}{c3-function-limit-arithmetic}
  Пусть $f\colon E\to\R$ и $g\colon E\to\R$ — две функции с общей
  областью определения. Если
  \[
    \lim_{E\ni x\to a}f(x)=A,
    \qquad
    \lim_{E\ni x\to a}g(x)=B,
  \]
  то
  \begin{enumerate}[label=\alph*)]
    \item $\displaystyle\lim_{E\ni x\to a}(f+g)(x)=A+B$;
    \item $\displaystyle\lim_{E\ni x\to a}(f\cdot g)(x)=AB$;
    \item $\displaystyle\lim_{E\ni x\to a}\left(\frac fg\right)(x)=\frac AB$,
      если $B\ne0$ и $g(x)\ne0$ при $x\in E$.
  \end{enumerate}
\end{theorem}

\begin{definition}{Бесконечно малая функция}{c3-infinitesimal-function}
  Функцию $f\colon E\to\R$ называют бесконечно малой при $E\ni x\to a$,
  если
  \[
    \lim_{E\ni x\to a}f(x)=0.
  \]
\end{definition}

\begin{proposition}{Свойства бесконечно малых функций}{c3-infinitesimal-properties}
  \begin{enumerate}[label=\alph*)]
    \item Если $\alpha\colon E\to\R$ и $\beta\colon E\to\R$ — бесконечно
      малые функции при $E\ni x\to a$, то их сумма
      $\alpha+\beta\colon E\to\R$ — также бесконечно малая функция при
      $E\ni x\to a$.
    \item Если $\alpha\colon E\to\R$ и $\beta\colon E\to\R$ — бесконечно
      малые функции при $E\ni x\to a$, то их произведение
      $\alpha\beta\colon E\to\R$ — также бесконечно малая функция при
      $E\ni x\to a$.
    \item Если $\alpha\colon E\to\R$ — бесконечно малая функция при
      $E\ni x\to a$, а $\beta\colon E\to\R$ — финально ограниченная
      функция при $E\ni x\to a$, то произведение
      $\alpha\beta\colon E\to\R$ есть бесконечно малая функция при
      $E\ni x\to a$.
  \end{enumerate}
\end{proposition}

\begin{theorem}{Предельный переход и неравенства}{c3-function-limit-inequalities}
  \begin{enumerate}[label=\alph*)]
    \item Если функции $f\colon E\to\R$ и $g\colon E\to\R$ таковы, что
      \[
        \lim_{E\ni x\to a}f(x)=A,
        \qquad
        \lim_{E\ni x\to a}g(x)=B
      \]
      и $A<B$, то найдется проколотая окрестность $\mathring U_E(a)$ точки
      $a$ в множестве $E$, в любой точке которой выполнено неравенство
      $f(x)<g(x)$.
    \item Если между функциями $f\colon E\to\R$, $g\colon E\to\R$ и
      $h\colon E\to\R$ на множестве $E$ имеет место соотношение
      $f(x)\le g(x)\le h(x)$ и если
      \[
        \lim_{E\ni x\to a}f(x)=
        \lim_{E\ni x\to a}h(x)=C,
      \]
      то существует также предел $g(x)$ при $E\ni x\to a$, причем
      $\lim\limits_{E\ni x\to a}g(x)=C$.
  \end{enumerate}
\end{theorem}

\begin{corollary}{Предельный переход в неравенствах}{c3-function-inequality-corollary}
  Пусть
  \[
    \lim_{E\ni x\to a}f(x)=A,
    \qquad
    \lim_{E\ni x\to a}g(x)=B.
  \]
  Если в некоторой проколотой окрестности $\mathring U_E(a)$ точки $a$
  \begin{enumerate}[label=\alph*)]
    \item выполнено $f(x)>g(x)$, то $A\ge B$;
    \item выполнено $f(x)\ge g(x)$, то $A\ge B$;
    \item выполнено $f(x)>B$, то $A\ge B$;
    \item выполнено $f(x)\ge B$, то $A\ge B$.
  \end{enumerate}
\end{corollary}

\begin{properties}{Показательная функция}{c3-exponential-properties}
  При $a>0$, $a\ne1$ функция $x\mapsto a^x$ обладает следующими
  свойствами:
  \begin{enumerate}[label=\arabic*)]
    \item $a^1=a$;
    \item $a^{x_1}a^{x_2}=a^{x_1+x_2}$;
    \item $a^x\to a^{x_0}$ при $x\to x_0$;
    \item $(a^{x_1}<a^{x_2})\Leftrightarrow(x_1<x_2)$, если $a>1$,
      и $(a^{x_1}>a^{x_2})\Leftrightarrow(x_1<x_2)$, если $0<a<1$;
    \item множеством значений функции $x\mapsto a^x$ является множество
      $\R_+=\{y\in\R\mid 0<y\}$ всех положительных чисел.
  \end{enumerate}
\end{properties}

\begin{definition}{Показательная функция}{c3-exponential-function}
  Отображение $x\mapsto a^x$ называется показательной или экспоненциальной
  функцией при основании $a$. Особенно часто встречается функция
  $x\mapsto e^x$, когда $a=e$, которую нередко обозначают через $\exp x$.
  В связи с этим для обозначения функции $x\mapsto a^x$ также иногда
  используется символ $\exp_a x$.
\end{definition}

\begin{definition}{Логарифмическая функция}{c3-logarithmic-function}
  Отображение, обратное к $\exp_a\colon\R\to\R_+$, называется
  логарифмической функцией при основании $a$ $(0<a$, $a\ne1)$ и
  обозначается символом
  \[
    \log_a\colon\R_+\to\R.
  \]
\end{definition}

\begin{definition}{Натуральный логарифм}{c3-natural-logarithm}
  При основании $a=e$ логарифмическая функция, или логарифм, называется
  натуральным логарифмом и обозначается $\ln\colon\R_+\to\R$.
\end{definition}

\begin{properties}{Логарифмическая функция}{c3-logarithm-properties}
  В области $\R_+$ своего определения логарифм обладает следующими
  свойствами:
  \begin{enumerate}[label=\arabic*${}^{\circ}$.]
    \item $\log_a a=1$;
    \item $\log_a(y_1y_2)=\log_a y_1+\log_a y_2$;
    \item $\log_a y\to\log_a y_0$ при $\R_+\ni y\to y_0\in\R_+$;
    \item $(\log_a y_1<\log_a y_2)\Leftrightarrow(y_1<y_2)$, если $a>1$,
      и $(\log_a y_1>\log_a y_2)\Leftrightarrow(y_1<y_2)$, если $0<a<1$;
    \item множество значений функции $\log_a\colon\R_+\to\R$ совпадает с
      множеством $\R$ всех действительных чисел;
    \item для любого $b>0$ и любого $\alpha\in\R$
      \[
        \log_a(b^\alpha)=\alpha\log_a b.
      \]
  \end{enumerate}
\end{properties}

\begin{properties}{Правило возведения степени в степень}{c3-power-of-power}
  Для любых $\alpha,\beta\in\R$ и $a>0$ имеет место равенство
  \[
    (a^\alpha)^\beta=a^{\alpha\beta}.
  \]
\end{properties}

\begin{definition}{Степенная функция}{c3-power-function}
  Функция $x\mapsto x^\alpha$, определенная на множестве $\R_+$
  положительных чисел, называется степенной функцией, а число $\alpha$
  называется показателем степени.
\end{definition}

\subsection{Общее определение предела функции (предел по базе)}

\begin{definition}{База}{c3-base}
  Совокупность $\mathcal B$ подмножеств $B\subset X$ множества $X$ будем
  называть базой в множестве $X$, если выполнены два условия:
  \[
    \mathrm{B}_1)\quad \forall B\in\mathcal B\;(B\ne\varnothing),
  \]
  \[
    \mathrm{B}_2)\quad
    \forall B_1\in\mathcal B\ \forall B_2\in\mathcal B\
    \exists B\in\mathcal B\;(B\subset B_1\cap B_2).
  \]
  Иными словами, элементы совокупности $\mathcal B$ суть непустые множества
  и в пересечении любых двух из них содержится некоторый элемент из той же
  совокупности.
\end{definition}

\begin{properties}{Наиболее употребительные базы}{c3-standard-bases}
  \begin{itemize}
    \item Запись $x\to a$ означает базу проколотых окрестностей точки
      $a\in\R$, элементами которой служат множества
      \[
        \mathring U(a):=
        \{x\in\R\mid a-\delta_1<x<a+\delta_2\ \land\ x\ne a\},
        \qquad \delta_1>0,\ \delta_2>0.
      \]
    \item Запись $x\to\infty$ означает базу окрестностей бесконечности,
      элементами которой служат множества
      \[
        U(\infty):=\{x\in\R\mid \delta<|x|\},
        \qquad \delta\in\R.
      \]
    \item Записи $x\to a$, $x\in E$, или $E\ni x\to a$ означают базу
      проколотых окрестностей точки $a$ в множестве $E$:
      \[
        \mathring U_E(a):=E\cap\mathring U(a),
      \]
      где предполагается, что $a$ — предельная точка множества $E$.
    \item Записи $x\to\infty$, $x\in E$, или $E\ni x\to\infty$ означают
      базу окрестностей бесконечности в множестве $E$:
      \[
        U_E(\infty):=E\cap U(\infty),
      \]
      где предполагается, что множество $E$ не ограничено.
  \end{itemize}

  Если $E=E_a^+=\{x\in\R\mid x>a\}$
  (соответственно $E=E_a^-=\{x\in\R\mid x<a\}$), то вместо $x\to a$,
  $x\in E$ пишут $x\to a+0$ (соответственно $x\to a-0$) и говорят, что
  $x$ стремится к $a$ справа или со стороны больших значений
  (соответственно слева или со стороны меньших значений).

  Если $E=E_\infty^+=\{x\in\R\mid c<x\}$
  (соответственно $E=E_\infty^-=\{x\in\R\mid x<c\}$), то вместо
  $x\to\infty$, $x\in E$ пишут $x\to+\infty$ (соответственно
  $x\to-\infty$) и говорят, что $x$ стремится к плюс бесконечности
  (соответственно к минус бесконечности).
\end{properties}

\begin{definition}{Предел функции по базе}{c3-limit-by-base}
  Пусть $f\colon X\to\R$ — функция на множестве $X$; $\mathcal B$ —
  база в $X$. Число $A\in\R$ называется пределом функции
  $f\colon X\to\R$ по базе $\mathcal B$, если для любой окрестности $V(A)$
  точки $A$ найдется элемент $B\in\mathcal B$ базы, образ которого $f(B)$
  содержится в окрестности $V(A)$.

  Если $A$ — предел функции $f\colon X\to\R$ по базе $\mathcal B$, то
  пишут
  \[
    \lim_{\mathcal B}f(x)=A.
  \]
  В логической символике:
  \[
    \lim_{\mathcal B}f(x)=A
    \;:=\;
    \forall V(A)\ \exists B\in\mathcal B\;(f(B)\subset V(A)),
  \]
  или, что равносильно,
  \[
    \lim_{\mathcal B}f(x)=A
    \;:=\;
    \forall\varepsilon>0\ \exists B\in\mathcal B\ \forall x\in B\
    (|f(x)-A|<\varepsilon).
  \]
\end{definition}

\begin{definition}{Бесконечные пределы по базе}{c3-infinite-limits-by-base}
  В соответствии с общим определением предела принимаются следующие
  соглашения:
  \[
    \lim_{\mathcal B}f(x)=\infty
    \;:=\;
    \forall V(\infty)\ \exists B\in\mathcal B\
    \bigl(f(B)\subset V(\infty)\bigr),
  \]
  или, что то же самое,
  \[
    \lim_{\mathcal B}f(x)=\infty
    \;:=\;
    \forall\varepsilon>0\ \exists B\in\mathcal B\ \forall x\in B\
    (\varepsilon<|f(x)|),
  \]
  \[
    \lim_{\mathcal B}f(x)=+\infty
    \;:=\;
    \forall\varepsilon\in\R\ \exists B\in\mathcal B\ \forall x\in B\
    (\varepsilon<f(x)),
  \]
  \[
    \lim_{\mathcal B}f(x)=-\infty
    \;:=\;
    \forall\varepsilon\in\R\ \exists B\in\mathcal B\ \forall x\in B\
    (f(x)<\varepsilon).
  \]
\end{definition}

\begin{definition}{Финально постоянная функция при базе}{c3-eventually-constant-base}
  Функция $f\colon X\to\R$ называется финально постоянной при базе
  $\mathcal B$, если существуют число $A\in\R$ и такой элемент
  $B\in\mathcal B$ базы, в любой точке $x\in B$ которого $f(x)=A$.
\end{definition}

\begin{definition}{Финально ограниченная функция при базе}{c3-eventually-bounded-base}
  Функция $f\colon X\to\R$ называется ограниченной при базе $\mathcal B$
  или финально ограниченной при базе $\mathcal B$, если существуют число
  $c\in\R$ и такой элемент $B\in\mathcal B$ базы, в любой точке $x\in B$
  которого $|f(x)|<c$.
\end{definition}

\begin{definition}{Бесконечно малая функция при базе}{c3-infinitesimal-base}
  Функция $f\colon X\to\R$ называется бесконечно малой при базе
  $\mathcal B$, если
  \[
    \lim_{\mathcal B}f(x)=0.
  \]
\end{definition}

\begin{properties}{Свойства предела по произвольной базе}{c3-arbitrary-base-limit-properties}
  Все свойства предела, установленные для специальной базы
  $E\ni x\to a$, справедливы для пределов по любой базе.
\end{properties}

\clearpage
\subsection{Вопросы существования предела функции}

\begin{definition}{Колебание функции}{c3-oscillation}
  Колебанием функции $f\colon X\to\R$ на множестве $E\subset X$
  называется величина
  \[
    \omega(f;E):=\sup_{x_1,x_2\in E}|f(x_1)-f(x_2)|,
  \]
  то есть верхняя грань модуля разности значений функции на всевозможных
  парах точек $x_1,x_2\in E$.
\end{definition}

\begin{theorem}{Критерий Коши существования предела функции}{c3-cauchy-criterion-function}
  Пусть $X$ — множество и $\mathcal B$ — база в $X$.

  Функция $f\colon X\to\R$ имеет предел по базе $\mathcal B$ в том и только
  в том случае, когда для любого числа $\varepsilon>0$ найдется элемент
  $B\in\mathcal B$ базы, на котором колебание функции меньше
  $\varepsilon$:
  \[
    \exists\lim_{\mathcal B}f(x)
    \quad\Longleftrightarrow\quad
    \forall\varepsilon>0\ \exists B\in\mathcal B\
    \bigl(\omega(f;B)<\varepsilon\bigr).
  \]
\end{theorem}

\begin{theorem}{О пределе композиции функций}{c3-composition-limit}
  Пусть $Y$ — множество; $\mathcal B_Y$ — база в $Y$;
  $g\colon Y\to\R$ — отображение, имеющее предел по базе
  $\mathcal B_Y$.

  Пусть $X$ — множество, $\mathcal B_X$ — база в $X$ и
  $f\colon X\to Y$ — такое отображение $X$ в $Y$, что для любого элемента
  $B_Y\in\mathcal B_Y$ базы $\mathcal B_Y$ найдется элемент
  $B_X\in\mathcal B_X$ базы $\mathcal B_X$, образ которого $f(B_X)$
  содержится в $B_Y$.

  При этих условиях композиция $g\circ f\colon X\to\R$ отображений $f$ и
  $g$ определена, имеет предел по базе $\mathcal B_X$ и
  \[
    \lim_{\mathcal B_X}(g\circ f)(x)=\lim_{\mathcal B_Y}g(y).
  \]
\end{theorem}

\begin{definition}{Монотонная функция}{c3-monotone-function}
  Функция $f\colon E\to\R$, определенная на числовом множестве
  $E\subset\R$, называется
  \begin{itemize}
    \item возрастающей на $E$, если
      \[
        \forall x_1,x_2\in E\;(x_1<x_2\Rightarrow f(x_1)<f(x_2));
      \]
    \item неубывающей на $E$, если
      \[
        \forall x_1,x_2\in E\;(x_1<x_2\Rightarrow f(x_1)\le f(x_2));
      \]
    \item невозрастающей на $E$, если
      \[
        \forall x_1,x_2\in E\;(x_1<x_2\Rightarrow f(x_1)\ge f(x_2));
      \]
    \item убывающей на $E$, если
      \[
        \forall x_1,x_2\in E\;(x_1<x_2\Rightarrow f(x_1)>f(x_2)).
      \]
  \end{itemize}
  Функции перечисленных типов называются монотонными на множестве $E$.
\end{definition}

\begin{theorem}{Критерий существования предела монотонной функции}{c3-monotone-function-limit}
  Предположим, что числа (или символы $-\infty$, $+\infty$)
  $i=\inf E$ и $s=\sup E$ являются предельными точками множества $E$ и
  $f\colon E\to\R$ — монотонная функция на $E$.

  Для того чтобы неубывающая на множестве $E$ функция $f\colon E\to\R$
  имела предел при $x\to s$, $x\in E$, необходимо и достаточно, чтобы она
  была ограничена сверху, а для того чтобы она имела предел при $x\to i$,
  $x\in E$, необходимо и достаточно, чтобы она была ограничена снизу.
\end{theorem}

\begin{definition}{Финальное свойство при базе}{c3-final-property}
  Условимся говорить, что некоторое свойство функций или соотношение между
  функциями выполнено финально при данной базе $\mathcal B$, если найдется
  элемент $B\in\mathcal B$ базы, на котором оно имеет место.
\end{definition}

\begin{definition}{Символ $o$ малое}{c3-little-o}
  Говорят, что функция $f$ есть бесконечно малая по сравнению с функцией
  $g$ при базе $\mathcal B$ и пишут $f=o(g)$ или
  $f=o_{\mathcal B}(g)$, если финально при базе $\mathcal B$ выполнено
  соотношение
  \[
    f(x)=\alpha(x)g(x),
  \]
  где $\alpha$ — функция, бесконечно малая при базе $\mathcal B$.
\end{definition}

\begin{definition}{Бесконечно малая более высокого порядка}{c3-higher-order-infinitesimal}
  Если $f=o_{\mathcal B}(g)$ и функция $g$ сама есть бесконечно малая при
  базе $\mathcal B$, то говорят, что $f$ есть бесконечно малая более высокого
  по сравнению с $g$ порядка при базе $\mathcal B$.
\end{definition}

\begin{definition}{Бесконечно большая функция}{c3-infinitely-large-function}
  Функцию, стремящуюся к бесконечности при данной базе, называют бесконечно
  большой функцией или просто бесконечно большой при данной базе.
\end{definition}

\begin{definition}{Бесконечно большая более высокого порядка}{c3-higher-order-infinite}
  Если $f$ и $g$ — бесконечно большие при базе $\mathcal B$ и
  $f=o_{\mathcal B}(g)$, то говорят, что $g$ есть бесконечно большая более
  высокого порядка по сравнению с $f$.
\end{definition}

\begin{definition}{Символ $O$ большое}{c3-big-o}
  Условимся, что запись $f=O(g)$ или $f=O_{\mathcal B}(g)$ при базе
  $\mathcal B$ будет означать, что финально при базе $\mathcal B$ выполнено
  соотношение
  \[
    f(x)=\beta(x)g(x),
  \]
  где $\beta(x)$ — финально ограниченная при базе $\mathcal B$ функция.

  В частности, запись $f=O_{\mathcal B}(1)$ означает, что функция $f$
  финально ограничена при базе $\mathcal B$.
\end{definition}

\begin{definition}{Функции одного порядка}{c3-same-order}
  Говорят, что функции $f$ и $g$ одного порядка при базе $\mathcal B$ и
  пишут $f\asymp_{\mathcal B}g$, если одновременно
  \[
    f=O_{\mathcal B}(g)
    \qquad\text{и}\qquad
    g=O_{\mathcal B}(f).
  \]

  Это условие равносильно тому, что найдутся числа $c_1>0$, $c_2>0$ и
  элемент $B\in\mathcal B$ такие, что на $B$ имеют место соотношения
  \[
    c_1|g(x)|\le |f(x)|\le c_2|g(x)|.
  \]
\end{definition}

\begin{definition}{Эквивалентные функции}{c3-equivalent-functions}
  Если между функциями $f$ и $g$ финально при базе $\mathcal B$ выполнено
  соотношение
  \[
    f(x)=\gamma(x)g(x),
    \qquad
    \lim_{\mathcal B}\gamma(x)=1,
  \]
  то говорят, что при базе $\mathcal B$ функция $f$ асимптотически ведет
  себя как функция $g$ или, короче, что $f$ эквивалентна $g$ при базе
  $\mathcal B$.

  В этом случае пишут $f\sim g$ или $f\sim_{\mathcal B}g$.
\end{definition}

\begin{properties}{Эквивалентность функций}{c3-equivalence-properties}
  \[
    f\sim_{\mathcal B}f,
  \]
  \[
    f\sim_{\mathcal B}g\quad\Longrightarrow\quad g\sim_{\mathcal B}f,
  \]
  \[
    f\sim_{\mathcal B}g\ \land\ g\sim_{\mathcal B}h
    \quad\Longrightarrow\quad
    f\sim_{\mathcal B}h.
  \]
  Кроме того, соотношение $f\sim_{\mathcal B}g$ равносильно тому, что
  \[
    f(x)=g(x)+o_{\mathcal B}(g(x)).
  \]
\end{properties}

\begin{proposition}{Замена на эквивалентную функцию в произведении}{c3-equivalent-substitution}
  Если $f\sim_{\mathcal B}\widetilde f$, то
  \[
    \lim_{\mathcal B}f(x)g(x)
    =
    \lim_{\mathcal B}\widetilde f(x)g(x),
  \]
  если один из этих пределов существует.
\end{proposition}

\begin{proposition}{Правила обращения с символами $o(\cdot)$ и $O(\cdot)$}{c3-o-O-rules}
  При данной базе:
  \begin{enumerate}[label=\alph*)]
    \item $o(f)+o(f)=o(f)$;
    \item $o(f)$ есть также $O(f)$;
    \item $o(f)+O(f)=O(f)$;
    \item $O(f)+O(f)=O(f)$;
    \item если $g(x)\ne0$, то
      \[
        \frac{o(f(x))}{g(x)}=o\!\left(\frac{f(x)}{g(x)}\right),
        \qquad
        \frac{O(f(x))}{g(x)}=O\!\left(\frac{f(x)}{g(x)}\right).
      \]
  \end{enumerate}
\end{proposition}
